\documentclass[11pt]{article}
\usepackage[a4paper,margin=26mm]{geometry}
\usepackage{amsmath,amssymb,amsthm,hyperref}
\newtheorem{theorem}{Theorem}
\newtheorem{lemma}[theorem]{Lemma}
\newtheorem{corollary}[theorem]{Corollary}
\title{Integer quadrature rounding gives $C(n!)^2$ fair dice}
\author{Research notes, 30 September 2026}
\date{}
\begin{document}
\maketitle
\noindent\textbf{Status.} Complete proof with independent mathematical reviews completed
by two collaborating agents. These reviews are not machine-checked formal
verification. The asymptotic class remains $\exp(O(n\log n))$, but the
leading coefficient in the exponent matches the whole-copy insertion
barrier proved in the companion note.

\begin{theorem}\label{thm:main}
There is an absolute positive integer $C$ such that, for every $n\ge1$,
there are $n$ equal-sided permutation-fair dice with
\[
 F_n=C(n!)^2
\]
faces each. One explicit admissible choice is $C=2^{25}$.
The construction adds one letter at a time, using $(r+1)^2$ unchanged
copies in the step from $r$ letters to $r+1$ letters.
\end{theorem}

The improvement comes from replacing common-denominator clearing of the
specific Hahn weights by rounding and exact integer corrections in a
unimodular binomial basis. The existing positivity estimates are retained.

\begin{lemma}[Integer quadrature rounding]\label{lem:round}
Let $N\ge r\ge1$, and suppose $w_0,\ldots,w_N$ form an exact degree-$r$
quadrature for the uniform probability measure on $[0,1]$, supported at
$i/N$, with $w_i\ge\eta>0$. Define
\[
 \mu_j=\int_0^1\binom{Nx}{j}\,dx\qquad(0\le j\le r).
\]
Suppose $K$ is a positive integer such that all $K\mu_j$ are integers.
If
\[
 K\eta\ge1+\max_{0\le i\le r}
       \sum_{j=i}^r\binom ji\binom{N+1}{j+1},
\]
then there are nonnegative integers $k_0,\ldots,k_N$ with
$\sum_i k_i=K$ whose normalized weights $k_i/K$ give another exact
degree-$r$ quadrature on this grid.
\end{lemma}
\begin{proof}
Set $a_i=\lfloor Kw_i\rfloor$ and
\[
 b_j=K\mu_j-\sum_{i=0}^Na_i\binom ij.
\]
The $b_j$ are integers. Exactness and $0\le Kw_i-a_i<1$ give
\[
 0\le b_j\le\sum_{i=0}^N\binom ij=\binom{N+1}{j+1}.
\]
The transposed Pascal matrix on nodes $0,\ldots,r$ is unimodular.
Its inverse gives integer corrections
\[
 d_i=\sum_{j=i}^r(-1)^{j-i}\binom ji b_j\qquad(0\le i\le r),
\]
which satisfy $\sum_{i=0}^r d_i\binom ij=b_j$ for every $j\le r$.
Put $k_i=a_i+d_i$ for $i\le r$, and $k_i=a_i$ otherwise.
The assumed bound and $a_i\ge K\eta-1$ imply $k_i\ge0$.
The corrected binomial moments are exact; these polynomials span the
degree-$r$ space. The zeroth moment gives $\sum_i k_i=K$.
\end{proof}

\begin{lemma}[The required moment denominators]
For integer $N$ and $0\le j\le r$, the denominator of $\mu_j$ divides
\[
 r!\operatorname{lcm}(1,2,\ldots,r+1).
\]
In particular $((r+1)!)^2\mu_j$ is an integer.
\end{lemma}
\begin{proof}
$j!\binom{Nx}{j}$ is a polynomial in $x$ with integer coefficients.
Integrating its monomials introduces only denominators $1,\ldots,j+1$.
Thus the denominator divides $j!\operatorname{lcm}(1,\ldots,j+1)$,
which divides the displayed quantity and then $((r+1)!)^2$.
\end{proof}

\begin{lemma}[A uniform positive starting quadrature]\label{lem:positive}
For every $r\ge1$ and $N=(r+1)^2$, the Wilson--Hahn projection weights
give an exact degree-$r$ uniform quadrature on $i/N$ satisfying
\[
 w_i>\frac1{8N}\qquad(0\le i\le N).
\]
\end{lemma}
\begin{proof}
Use the manuscript's Hahn polynomials with largest grid index $N$.
Its estimates apply for $N\ge r(r+2)$, and show, for $r\ge2$, that the
unnormalized weights satisfy
\[
 W_i>\frac{N-3r(r+3)/4}{N+1}>\frac18.
\]
The last inequality follows from
$7(r+1)^2-6r(r+3)-1=(r-2)^2+2>0$.
For $r=1$, all $W_i=N/(N+1)>1/8$ directly. Normalize by $N$.
The applicability of the manuscript estimates uses
$r(r+1)\le N$ for Zaremba's nodal bound and
$(r+1)^2\le N+1$ for Wilson's integral estimate.
\end{proof}

\begin{lemma}[An absolute scaling constant suffices]\label{lem:constant}
For every integer $s\ge2$,
\[
 \frac{8s^2\left(2^s\binom{s^2+1}{s}+1\right)}{(s!)^2}<2^{25}.
\]
\end{lemma}
\begin{proof}
The second summand, $8s^2/(s!)^2$, is at most $8$.
For the first, use
\[
 \binom{s^2+1}{s}\le\frac{s^{2s}e^{1/s}}{s!},
 \qquad s!\ge\sqrt{2\pi s}(s/e)^s.
\]
It is consequently at most
\[
 \frac{8e^{1/s}}{(2\pi)^{3/2}}\sqrt{s}
       \left(\frac{2e^3}{s}\right)^s
 <\sqrt{s}\left(\frac{81/2}{s}\right)^s.
\]
Here $\pi>3$ and $e<2.72$ justify both strict inequalities, uniformly
for $s\ge2$. Since $\sqrt{s}<2e^{s/16}$ and
$\sup_{u>0}(a/u)^u=e^{a/e}$, this is bounded by
\[
 2\exp\left(\frac{81}{2}e^{-15/16}\right)<2e^{16}.
\]
The inequality $e^{15/16}>81/32$ follows by retaining its first five
Taylor terms. Finally
$2e^{16}+8<2(11/4)^{16}+8<2^{25}$.
\end{proof}

\begin{proof}[Proof of Theorem~\ref{thm:main}]
Start with the one-letter string $s_1=1^C$, where $C=2^{25}$.
Suppose $s_r$ is $r$-fair and every letter occurs $C(r!)^2$ times.
Set
\[
 N=(r+1)^2,\qquad K=C((r+1)!)^2.
\]
The old letters each occur $NC(r!)^2=K$ times in $N$ copies of $s_r$.
The moment-denominator lemma says $K\mu_j$ is integral for every $j\le r$.
Use the starting quadrature from Lemma~\ref{lem:positive}.
Because $N=(r+1)^2$, the binomial coefficients
$\binom{N+1}{j+1}$ increase for $0\le j\le r$, and hence
\[
 \sum_{j=i}^r\binom ji\binom{N+1}{j+1}
 \le\binom{N+1}{r+1}\binom{r+1}{i+1}
 \le2^{r+1}\binom{N+1}{r+1}.
\]
Lemma~\ref{lem:constant}, with $s=r+1$, verifies the sufficient
positivity condition in Lemma~\ref{lem:round}, using $\eta=1/(8N)$.
Obtain integer gap lengths $k_i$ of total $K$ and exact degree-$r$
Bernstein moments. The insertion criterion proves that the resulting
string $s_{r+1}$ is permutation-fair. Its new and old letters all occur
$K$ times, completing the induction.
\end{proof}

\begin{corollary}[Optimal leading exponent within whole-copy insertion]
The constructed dice satisfy
\[
 \log F_n=2n\log n-2n+O(\log n).
\]
The companion insertion-barrier note proves that every one-letter-at-a-time
whole-copy insertion construction has an oldest die with
$\log F_n\ge2n\log n-O(n)$. Thus the leading coefficient $2$ is
optimal for that framework.
\end{corollary}

The constant is deliberately conservative. An exact finite check of the
displayed sufficient ratio for $2\le s\le128$ gives a maximum below
$2^{22}$ (near $s=15$); an additional tail estimate would certify that
smaller constant. This optimization has no bearing on the asymptotic result.

\end{document}
